\documentclass[11pt,a4paper]{article}
\usepackage[margin=25mm]{geometry}
\usepackage{amsmath,amssymb,amsthm,booktabs,hyperref}
\setlength{\emergencystretch}{2em}
\newtheorem{lemma}{Lemma}
\newtheorem{theorem}[lemma]{Theorem}
\newcommand{\gen}[1]{\langle#1\rangle}
\newcommand{\Sym}{\operatorname{Sym}}
\newcommand{\Alt}{\operatorname{Alt}}
\title{An even attaching group and the critical symmetric wing}
\author{Human Compression project: a component of the L019 recovery}
\date{}
\begin{document}\maketitle
An endpoint anchor of an ordered involution family is a point moved by
its first member and fixed by every later member. All groups below are
actual permutation groups. An attachment from outside its domain at
an anchor $a$ has the form $(a\ c)\gamma$, where $c$ is new and
$\gamma$ is an involution or identity supported off the old domain
and $c$. In particular, $\gamma$ centralizes every old mark and the
pure transposition $(a\ c)$.

\begin{lemma}\label{ten}
Let $(t_0,\ldots,t_4)$ be a natural rank-five C-marking of $A_{10}$.
Suppose $a$ is an endpoint anchor and its tail is transitive on the
other nine points. No attachment $(a\ c)\gamma$ as above can precede
this marking in a string C-marking.
\end{lemma}
\begin{proof}
We acquire the complete ordered rank-five list from the accepted
small-degree part of the alternating rank theorem (L044). Its retained terminal list has ten rank-five tuples: eight generate
$A_{10}$, while the other two generate groups of order $1920$.
We retain and scan this complete superset. The completeness of its
$A_{10}$ portion, including the original census validation, is an
acquired classification premise; the following operation only filters
that list.

For each retained tuple in each orientation, inspect every point
for the anchor condition, then inspect the tail orbits. Exactly four
oriented records pass, all belonging to the eight $A_{10}$ tuples. The table uses zero-based source indices and
orientation $0$ for the recorded order, $1$ for reversal. Set
\[
 u=(a\ 11),\quad U=\gen{t_0,t_1,t_2,t_3},\quad
 V=\gen{u,t_0,t_1,t_2},\quad M=\gen{t_0,t_1,t_2}.
\]
\begin{center}
\begin{tabular}{ccc rrrr}
\toprule
Index&Orientation&Anchor&$|U|$&$|V|$&$|U\cap V|$&$|M|$\\
\midrule
1&0&10&5040&5040&360&60\\
2&1&1&5040&5040&360&60\\
4&1&1&40320&10080&720&120\\
7&0&10&40320&10080&720&120\\
\bottomrule
\end{tabular}
\end{center}
The certificate acquires all ten literal tuples from the retained
list, rather than selecting only these four examples. The separate
GAP check identifies their exact group orders as well. A direct
permutation implementation enumerates each of the displayed finite
closures and intersects their elements; a separate GAP calculation
agrees. Every row gives $[U\cap V:M]=6$.

For an actual attachment $u\gamma$, the corresponding group on the
five entries $u\gamma,t_0,t_1,t_2,t_3$ lies in the direct product
of the pure group and $\gen\gamma$. Projection onto the pure group
has central kernel of order at most two. The acquired central
intersection lemma says that, if the actual marking had IP, the
index $[U\cap V:M]$ would divide a power of two: its difference
quotient embeds in a quotient of that central kernel. The factor
three in six contradicts this necessary condition. This includes
$\gamma=1$ and proves the assertion.
\end{proof}

\begin{theorem}\label{primitive}
Let $(t_0,\ldots,t_{k-1})$ be a faithful primitive string C-marking
on a set $Y$ of size $d\ge3$, consisting entirely of even
permutations. Suppose it has an endpoint anchor $a$, its tail is
transitive on $Y\setminus\{a\}$, and an attachment from outside
at $a$ occurs immediately before it in a larger string C-marking.
Then either
\[
                         2k\le d-1,                 \tag{1}
\]
or $d=6$, $k=3$, the group is the primitive degree-six copy of
$A_5$, and its tail has order ten.
\end{theorem}
\begin{proof}
The group is not $S_d$, because every mark is even. First suppose
it is the natural $A_d$. The complete acquired alternating rank
bound gives (1) for $d\ge12$. Its exact small-degree conclusions
leave only $A_5$ of rank three, $A_{10}$ of rank five, and $A_{11}$
of rank six as possible violations of (1).

The degree-five case has no endpoint anchor. Indeed a nonidentity
involution of $A_5$ has one fixed point. The third mark fixes $a$,
so its unique fixed point is $a$; the first mark commutes with it
and therefore fixes $a$, a contradiction.
The degree-ten case is excluded by Lemma~\ref{ten}, using the
explicit tail-transitivity hypothesis.
For degree eleven, the acquired actual-wing certificate inspects
all four complete rank-six tuples and both orientations. Each
nonempty anchor gives the same pure-attachment prefix index six,
with orders $5040,5040,360,60$. The same central-kernel argument
excludes any attachment $(a\ c)\gamma$. Here the attachment is
already pure on $Y\cup\{c\}$, so no separate centralizer
normalization or general quotient-IP assertion is needed.

For every other primitive group, the complete primitive theorem
acquired in L047 applies to the actual faithful C-marking. Unless
its rank is strictly less than $d/2$, it belongs to the twelve
retained ordered small representatives: degree-six copies of
$A_5$ or $S_5$, or the degree-ten rank-five copy of $S_6$.
The already acquired twenty-four-orientation filter records every
anchor and complementary tail orbit. Without imposing its
additional character test, the even groups with a transitive tail
are exactly source indices and orientations
\[
                  (0,1),\ (1,0),\ (1,1),\ (2,0).
\]
They have degree six, rank three, group order sixty and tail order
ten. The other primitive representatives either have no anchor
with a transitive complementary tail or have an odd mark. We
explicitly do not reuse the character condition from the earlier
faithful-augmentation theorem: the present hypotheses supply no
such character. If the exceptional list does not apply, $k<d/2$
implies (1) by integrality. This completes the inventory.
\end{proof}

\begin{theorem}\label{descent}
Let an internal perfect split in a transitive natural symmetric
C-marking of degree $n$ and rank $r$, with $2r>n+2$, be
\[
 q=\rho_i=(a\ b)\beta,\qquad
 L=\gen{\rho_0,\ldots,\rho_{i-1}},\quad
 R=\gen{\rho_{i+1},\ldots,\rho_{r-1}},
\]
where $L$ is transitive and supported on $A$, $R$ is transitive
and supported on $B$, and $A\sqcup B$ is the domain. Suppose
$J=\gen{q,R}$ contains only even permutations. Then
$L=\Sym(A)$ and either
\[
                         2i>|A|+2,                  \tag{2}
\]
or $|B|=5$, $R$ has rank two and order ten, $|A|$ is even,
and $i=|A|/2+1$.
\end{theorem}
\begin{proof}
The right wing is even and nontrivial, so is not its full symmetric
group. The acquired actual-wing theorem therefore gives
$L=\Sym(A)$. The group $J$ is supported faithfully on
$Y=\{a\}\cup B$. It is transitive there, and its stabilizer at
$a$ contains $R$, transitive on the complement. Thus $J$ is doubly
transitive, hence primitive. Its actual submarking inherits IP.
Its first mark $q$ moves $a$, and its tail fixes $a$.

The adjacent left mark $\rho_{i-1}$ moves $a$: all earlier left
marks fix $a$ by commutation with the unique crossing, whereas
$L$ is transitive. Write $\rho_{i-1}=(a\ c)\gamma$. Its support
is in $A$, so $c\in A\setminus\{a\}$ and $\gamma$ has support
disjoint from $Y\cup\{c\}$. This is exactly the actual outside
attachment required in Theorem~\ref{primitive}; no restriction
kernel is being silently discarded.

Put $t=\operatorname{rank}(R)$ and $b=|B|$.
In the ordinary case that theorem gives $2(t+1)\le b$.
Consequently
\[
  2i=2r-2-2t>|A|+b-2t\ge|A|+2,
\]
which is (2). Otherwise $b=5$ and $t=2$. Then
$2i>|A|+1$; if (2) fails, the even integer $2i$ equals
$|A|+2$, giving the stated critical case and parity of $|A|$.
\end{proof}

\paragraph{Scope and attribution.}
The alternating and transitive classification theorems, their exact
finite lists and their complete validation remain attributed to
Cameron, Fernandes, Leemans and Mixer through L044 and L047. The
A11 obstruction and central-intersection mechanism are acquired
from earlier separately reviewed L019 components. The new finite
work filters the complete A10 list and proves the changed even
attachment interface, with its precise residual degree-six case.
The intended high-rank extension problem is that of
Cameron--Fernandes--Leemans, arXiv:2212.12723v2.

This is a descent reduction, not an extension theorem for the
remaining critical case. In particular, a critical symmetric wing
need not admit a split with two odd attaching groups: an explicit
rank-seven $S_{12}$ marking can have only one odd mark, which is
pure. That diagnostic refutes a proposed strengthening of our own,
not a source conclusion. Valid leaf kernels, the full reversible
correspondence and the whole L019 principal claim remain pending.
No new census, open conjecture, Lean theorem or compression
measurement is claimed here.
\end{document}
